%% ---------------------------------------------------------------------------
%% House layout for Springer Nature / Nature-portfolio submissions.
%%
%%   %% ---------------------------------------------------------------------------
%% House layout for Springer Nature / Nature-portfolio submissions.
%%
%%   %% ---------------------------------------------------------------------------
%% House layout for Springer Nature / Nature-portfolio submissions.
%%
%%   %% ---------------------------------------------------------------------------
%% House layout for Springer Nature / Nature-portfolio submissions.
%%
%%   \input{sn-house-layout}        %% immediately after \documentclass
%%   ... \maketitle ...
%%   \linenumbers                   %% immediately after \maketitle
%%
%% Three settings, each with the measurement that chose it. Delete this file and
%% its two lines to return to the class's own layout; nothing else depends on it.
%% ---------------------------------------------------------------------------

%% 1. THE PAGE. The class loads geometry ITSELF, with text={31pc,194.25mm} on A4
%%    under sn-nature, so a second \usepackage{geometry} errors and the block has
%%    to be reset with \geometry instead.
%%
%%    A4 is 210 by 297 mm, so 170 by 257 leaves 20 mm on every side. Measured
%%    consequence rather than an aesthetic one: artwork authored at the venue's
%%    179 mm maximum and inserted at \textwidth is scaled by textwidth/179, which
%%    the class's own block makes 0.73 and this makes 0.950. A 7 pt label in the
%%    source therefore renders at about 5.11 pt under the class and 6.65 pt here,
%%    and 5.11 pt is below every legibility floor a venue states.
\geometry{textwidth=170mm, textheight=257mm, top=20mm, left=20mm,
          footskip=12mm, marginparsep=3mm, marginparwidth=8mm,
          bindingoffset=0mm}

%% 2. THE FOLIO. It is already bottom-centred by the class's \ps@headings and
%%    needs nothing. **Do not "fix" it with \pagestyle{plain}.** This class's
%%    plain style is
%%        \def\@oddfoot{\vbox to 18pt{...\hfil ddd\thepage\hfil}}
%%    which prints a literal "ddd" before every page number. Checked in
%%    sn-jnl.cls at line 870; the same bytes ship in the template and in every
%%    project copy taken from it.

%% 3. LINE NUMBERS. Use the lineno PACKAGE, not the class's own `lineno` option,
%%    and the reason is measured rather than stylistic.
%%
%%    The class option turns on `vruler` with
%%        \setvruler[12bp][1][1][3][1][1.18\textwidth][26pt][-7pt][0.99\textheight]
%%    which pads the counter to THREE digits and places it at 1.18\textwidth on
%%    odd pages. At the block above that lands the numerals at x = 203.0 mm on a
%%    210 mm page, measured from a compiled test, so a three-digit number clears
%%    the trim by about 0.7 mm and a fourth digit does not clear it at all. A
%%    manuscript of any length reaches four digits: the paper this file was
%%    extracted from runs to line 1500 over 28 pages.
%%
%%    The lineno package puts them in the left margin instead, measured at
%%    x = 13.7 mm against a text block starting at 20 mm, on every page and at
%%    every digit count. `mathlines` numbers display equations too. It must load
%%    before hyperref, which is why this file is input right after \documentclass.
\usepackage[mathlines]{lineno}

%% \linenumbers goes after \maketitle, not here, so the title block stays
%% unnumbered. Putting it in this file would number the authors and affiliations.
        %% immediately after \documentclass
%%   ... \maketitle ...
%%   \linenumbers                   %% immediately after \maketitle
%%
%% Three settings, each with the measurement that chose it. Delete this file and
%% its two lines to return to the class's own layout; nothing else depends on it.
%% ---------------------------------------------------------------------------

%% 1. THE PAGE. The class loads geometry ITSELF, with text={31pc,194.25mm} on A4
%%    under sn-nature, so a second \usepackage{geometry} errors and the block has
%%    to be reset with \geometry instead.
%%
%%    A4 is 210 by 297 mm, so 170 by 257 leaves 20 mm on every side. Measured
%%    consequence rather than an aesthetic one: artwork authored at the venue's
%%    179 mm maximum and inserted at \textwidth is scaled by textwidth/179, which
%%    the class's own block makes 0.73 and this makes 0.950. A 7 pt label in the
%%    source therefore renders at about 5.11 pt under the class and 6.65 pt here,
%%    and 5.11 pt is below every legibility floor a venue states.
\geometry{textwidth=170mm, textheight=257mm, top=20mm, left=20mm,
          footskip=12mm, marginparsep=3mm, marginparwidth=8mm,
          bindingoffset=0mm}

%% 2. THE FOLIO. It is already bottom-centred by the class's \ps@headings and
%%    needs nothing. **Do not "fix" it with \pagestyle{plain}.** This class's
%%    plain style is
%%        \def\@oddfoot{\vbox to 18pt{...\hfil ddd\thepage\hfil}}
%%    which prints a literal "ddd" before every page number. Checked in
%%    sn-jnl.cls at line 870; the same bytes ship in the template and in every
%%    project copy taken from it.

%% 3. LINE NUMBERS. Use the lineno PACKAGE, not the class's own `lineno` option,
%%    and the reason is measured rather than stylistic.
%%
%%    The class option turns on `vruler` with
%%        \setvruler[12bp][1][1][3][1][1.18\textwidth][26pt][-7pt][0.99\textheight]
%%    which pads the counter to THREE digits and places it at 1.18\textwidth on
%%    odd pages. At the block above that lands the numerals at x = 203.0 mm on a
%%    210 mm page, measured from a compiled test, so a three-digit number clears
%%    the trim by about 0.7 mm and a fourth digit does not clear it at all. A
%%    manuscript of any length reaches four digits: the paper this file was
%%    extracted from runs to line 1500 over 28 pages.
%%
%%    The lineno package puts them in the left margin instead, measured at
%%    x = 13.7 mm against a text block starting at 20 mm, on every page and at
%%    every digit count. `mathlines` numbers display equations too. It must load
%%    before hyperref, which is why this file is input right after \documentclass.
\usepackage[mathlines]{lineno}

%% \linenumbers goes after \maketitle, not here, so the title block stays
%% unnumbered. Putting it in this file would number the authors and affiliations.
        %% immediately after \documentclass
%%   ... \maketitle ...
%%   \linenumbers                   %% immediately after \maketitle
%%
%% Three settings, each with the measurement that chose it. Delete this file and
%% its two lines to return to the class's own layout; nothing else depends on it.
%% ---------------------------------------------------------------------------

%% 1. THE PAGE. The class loads geometry ITSELF, with text={31pc,194.25mm} on A4
%%    under sn-nature, so a second \usepackage{geometry} errors and the block has
%%    to be reset with \geometry instead.
%%
%%    A4 is 210 by 297 mm, so 170 by 257 leaves 20 mm on every side. Measured
%%    consequence rather than an aesthetic one: artwork authored at the venue's
%%    179 mm maximum and inserted at \textwidth is scaled by textwidth/179, which
%%    the class's own block makes 0.73 and this makes 0.950. A 7 pt label in the
%%    source therefore renders at about 5.11 pt under the class and 6.65 pt here,
%%    and 5.11 pt is below every legibility floor a venue states.
\geometry{textwidth=170mm, textheight=257mm, top=20mm, left=20mm,
          footskip=12mm, marginparsep=3mm, marginparwidth=8mm,
          bindingoffset=0mm}

%% 2. THE FOLIO. It is already bottom-centred by the class's \ps@headings and
%%    needs nothing. **Do not "fix" it with \pagestyle{plain}.** This class's
%%    plain style is
%%        \def\@oddfoot{\vbox to 18pt{...\hfil ddd\thepage\hfil}}
%%    which prints a literal "ddd" before every page number. Checked in
%%    sn-jnl.cls at line 870; the same bytes ship in the template and in every
%%    project copy taken from it.

%% 3. LINE NUMBERS. Use the lineno PACKAGE, not the class's own `lineno` option,
%%    and the reason is measured rather than stylistic.
%%
%%    The class option turns on `vruler` with
%%        \setvruler[12bp][1][1][3][1][1.18\textwidth][26pt][-7pt][0.99\textheight]
%%    which pads the counter to THREE digits and places it at 1.18\textwidth on
%%    odd pages. At the block above that lands the numerals at x = 203.0 mm on a
%%    210 mm page, measured from a compiled test, so a three-digit number clears
%%    the trim by about 0.7 mm and a fourth digit does not clear it at all. A
%%    manuscript of any length reaches four digits: the paper this file was
%%    extracted from runs to line 1500 over 28 pages.
%%
%%    The lineno package puts them in the left margin instead, measured at
%%    x = 13.7 mm against a text block starting at 20 mm, on every page and at
%%    every digit count. `mathlines` numbers display equations too. It must load
%%    before hyperref, which is why this file is input right after \documentclass.
\usepackage[mathlines]{lineno}

%% \linenumbers goes after \maketitle, not here, so the title block stays
%% unnumbered. Putting it in this file would number the authors and affiliations.
        %% immediately after \documentclass
%%   ... \maketitle ...
%%   \linenumbers                   %% immediately after \maketitle
%%
%% Three settings, each with the measurement that chose it. Delete this file and
%% its two lines to return to the class's own layout; nothing else depends on it.
%% ---------------------------------------------------------------------------

%% 1. THE PAGE. The class loads geometry ITSELF, with text={31pc,194.25mm} on A4
%%    under sn-nature, so a second \usepackage{geometry} errors and the block has
%%    to be reset with \geometry instead.
%%
%%    A4 is 210 by 297 mm, so 170 by 257 leaves 20 mm on every side. Measured
%%    consequence rather than an aesthetic one: artwork authored at the venue's
%%    179 mm maximum and inserted at \textwidth is scaled by textwidth/179, which
%%    the class's own block makes 0.73 and this makes 0.950. A 7 pt label in the
%%    source therefore renders at about 5.11 pt under the class and 6.65 pt here,
%%    and 5.11 pt is below every legibility floor a venue states.
\geometry{textwidth=170mm, textheight=257mm, top=20mm, left=20mm,
          footskip=12mm, marginparsep=3mm, marginparwidth=8mm,
          bindingoffset=0mm}

%% 2. THE FOLIO. It is already bottom-centred by the class's \ps@headings and
%%    needs nothing. **Do not "fix" it with \pagestyle{plain}.** This class's
%%    plain style is
%%        \def\@oddfoot{\vbox to 18pt{...\hfil ddd\thepage\hfil}}
%%    which prints a literal "ddd" before every page number. Checked in
%%    sn-jnl.cls at line 870; the same bytes ship in the template and in every
%%    project copy taken from it.

%% 3. LINE NUMBERS. Use the lineno PACKAGE, not the class's own `lineno` option,
%%    and the reason is measured rather than stylistic.
%%
%%    The class option turns on `vruler` with
%%        \setvruler[12bp][1][1][3][1][1.18\textwidth][26pt][-7pt][0.99\textheight]
%%    which pads the counter to THREE digits and places it at 1.18\textwidth on
%%    odd pages. At the block above that lands the numerals at x = 203.0 mm on a
%%    210 mm page, measured from a compiled test, so a three-digit number clears
%%    the trim by about 0.7 mm and a fourth digit does not clear it at all. A
%%    manuscript of any length reaches four digits: the paper this file was
%%    extracted from runs to line 1500 over 28 pages.
%%
%%    The lineno package puts them in the left margin instead, measured at
%%    x = 13.7 mm against a text block starting at 20 mm, on every page and at
%%    every digit count. `mathlines` numbers display equations too. It must load
%%    before hyperref, which is why this file is input right after \documentclass.
\usepackage[mathlines]{lineno}

%% \linenumbers goes after \maketitle, not here, so the title block stays
%% unnumbered. Putting it in this file would number the authors and affiliations.
